\documentclass[11pt,letterpaper]{article}
\usepackage[T1]{fontenc}
\usepackage[utf8]{inputenc}
\usepackage{lmodern,microtype}
\usepackage[margin=0.88in]{geometry}
\usepackage{amsmath,amssymb,amsthm,mathtools}
\usepackage{booktabs,longtable,tabularx,array}
\usepackage{xcolor,xurl,hyperref,enumitem,fancyhdr,needspace}
\hypersetup{colorlinks=true,linkcolor=blue!40!black,citecolor=blue!40!black,
 urlcolor=blue!45!black,pdftitle={PA194J and PA194D: two new 194-state PA consistency machines},
 pdfauthor={Research report for the ProveIt PA project}}
\urlstyle{tt}
\setlength{\emergencystretch}{3em}
\setlength{\parindent}{0pt}
\setlength{\parskip}{5pt}
\setlist{itemsep=3pt,topsep=5pt}
\pagestyle{fancy}\fancyhf{}
\fancyhead[L]{\small PA194J and PA194D: PA consistency machines}
\fancyhead[R]{\small 13 November 2026}
\fancyfoot[C]{\thepage}
\setlength{\headheight}{14pt}
\newcommand{\PA}{\mathrm{PA}}
\newcommand{\Con}{\operatorname{Con}}
\newcommand{\Halt}{\operatorname{Halt}}
\newcommand{\Prf}{\operatorname{Prf}}
\newcommand{\file}[1]{\texttt{\nolinkurl{#1}}}
\newcommand{\M}{\mathsf{PA194J}}
\newcommand{\Parent}{\mathsf{PA194CS}}
\newtheorem{theorem}{Theorem}[section]
\newtheorem{lemma}[theorem]{Lemma}
\newtheorem{proposition}[theorem]{Proposition}
\newtheorem{corollary}[theorem]{Corollary}
\title{PA194J and PA194D:\\two new 194-state PA consistency machines\\[0.45em]
\large Corpus comparison, zero-register jumps,\\a changed diagonal emitter, and PA-provable equivalence}
\author{Research report for the ProveIt PA project}
\date{13 November 2026, Pacific time}
\begin{document}
\maketitle
\begin{abstract}
A machine with fewer than 194 working states was not obtained. We provide
the requested alternative: a materially different 194-state, two-symbol,
single-tape machine, PA194J, with an ordinary mathematical proof of
\(\PA\vdash[\neg\Halt(\M,\mathrm{blank})\leftrightarrow\Con(\PA)]\).
The new machine replaces ten unconditional source jumps by decrements of
registers proved zero at those locations. A finite inductive abstraction
with 263,180 location/status pairs establishes exact equality of the
parent and child source runs from their common seed. Three transition-row
changes implement these substitutions, retaining the symmetric Shift
controller and the existing 67-reader decoder.

The construction contains 194 working states and 388 transition rows.
Its source has 3,646 primitive sites and 4,755 branch outcomes. A literal
tape checker derives four generic list lemmas from the actual rows and
checks the source outcomes for arbitrary natural register values, gap,
and finite nonempty spare list. Exact blank initialization takes 60,062
transitions and reaches the complete expected tape. The new graph differs
from every existing 194-state graph in the pinned corpus, even allowing
state renaming, a different designated start, reflection, and global
symbol exchange.

A second construction, PA194D, changes the logical emitter from
\(\Gamma(A)\) to \(\Gamma(D)\), with a correspondingly changed selector
in the strong proof-word compiler. A new allocation of 69 decoder readers
admits nine simultaneous history folds, giving \(134+69-9=194\) working
states. Its exact blank startup takes 58,168 transitions. It is supplied
as a separate literal table and construction, with its own calibration
argument and physical checks.

The comparison covers 41 research packets at ProveIt revision
\nolinkurl{17f92105e2bc4317feae6c219cb2143437b482ff}: 36 filed packets and
five newer incoming archives. The census recognizes 2,685 table
occurrences, 620 byte forms, and 446 action-graph classes; rejected and
diagnostic tables are identified separately. The proof transfers the
supplied ordinary PA calibration through new source and physical
preservation arguments. No new end-to-end Lean or Rocq theorem, iPRA
equivalence, state lower bound, or unconditional nonhalting is claimed.
\end{abstract}
\clearpage
\begingroup
\small
\setlength{\parskip}{0pt}
\tableofcontents
\endgroup
\clearpage

\section{Result and precise machine convention}

The principal file is \file{construction/machine/pa194j.tm}. The machine
starts in working state 0, scanning cell 0 of a two-way infinite tape
filled with zero. The alphabet is \(\{0,1\}\). Working states are
\(0,\ldots,193\); label 194 is the terminal halt control, has no outgoing
rows, and is excluded from the state count. Each transition writes a bit,
moves exactly one cell left or right, and enters its target state.
Every working state has exactly two rows. The five-column format is
\begin{center}
\file{state read_bit write_bit L_or_R next_state}.
\end{center}
All 194 working controls are reachable in the pointed transition graph.
This finite graph fact does not assert that the blank execution visits
every control.

The authoritative dense table has SHA-256
\begin{quote}\small
\nolinkurl{7c498ced0eb0dffed77ae2f68f4d80c98edfd8f7763ec1be68aa8cd7ea80eacd}.
\end{quote}
Its native-label presentation starts at 74 and halts at 254, with SHA-256
\begin{quote}\small
\nolinkurl{31abf450ba25066c8ac3aad6e7eedb59386b770bb4917f006800145ebb622218}.
\end{quote}
The supplied label map conjugates every transition exactly. The entire
dense table is printed in Appendix~\ref{app:table}.

\begin{theorem}[Main result]
For the literal table specified above, with \(H_\M\) the usual arithmetic
sentence asserting a finite halting trace from its blank start,
\[
 \PA\vdash \bigl(H_\M\leftrightarrow\neg\Con(\PA)\bigr),
 \qquad
 \PA\vdash \bigl(\neg H_\M\leftrightarrow\Con(\PA)\bigr).
\]
\end{theorem}

Here the proof is an ordinary mathematical proof formalizable in PA,
with executable finite and symbolic certificates. Its logical calibration
is inherited from the supplied PA194/PA195 developments~\cite{pa194,cached} and retained in
\file{references/logical_calibration/}. The machine-specific work in this
packet is the new source invariant, exact table construction, independent
physical simulation, and exact initialization. The repository's unfinished
iPRA endpoints and the absence of a new kernel export remain distinct
from this ordinary proof; see Section~\ref{sec:calibration}.

\subsection{Contribution and cost}

The immediate parent is the supplied PA194CS machine, which already uses
a symmetric destructive comparison in Shift. PA194J specializes its
control transfers using \emph{global source invariants}. At eight sites
register \(A\) is zero; at the scheduler site register \(X\) is zero; at
the final AX return register \(D\) is zero. At these ten sites a decrement
can implement precisely the former jump. These substitutions create new
pure right-moving readers in the physical decoder.

The count stays
\[
 \underbrace{135}_{\text{fixed part before the CS fold}}
 +\underbrace{67}_{\text{raw AX readers}}
 -\underbrace{8}_{\text{seven AX folds and the CS fold}}
 =194.
\]
Equivalently, after the controller fold there are 134 fixed controls and
60 AX controls. The new construction improves neither this count nor the
startup time: initialization grows from the parent's 57,768 transitions
to 60,062. Its value as the requested fallback is a different verified
physical implementation with new decoder interfaces and an explicit
global invariant available for further state reductions.

This is not merely a different numbering of an old table. The action
graphs are distinct under the broad finite equivalence tests in
Section~\ref{sec:novelty}. Two changed rows are exercised 3 and 21 times,
respectively, during the actual blank initialization. At source level the
program now has ten conditional instructions where the parent had jumps;
the proof explains why their zero branches always apply.

\section{Comparison of the supplied designs}
\label{sec:comparison}

\subsection{Scope and reproducible census}

All comparisons are pinned to the revision on the title page~\cite{repo}. The
\file{PA/docs} ledger contains 36 delivered packets. Five later incoming
archives are absent from that ledger and were read separately:
\begin{itemize}
\item \file{pa194C_conpa_2026-11-13 (2).zip};
\item \file{pa194_new_designs_2026-11-13.zip};
\item \file{pa194c_conpa.zip};
\item \file{pa194c_conpa_2026-11-13 (1).zip};
\item \file{pa194c_conpa_2026-11-13.zip}.
\end{itemize}
The byte identities of these five archives were checked against their
Git blob identities. All five repeat the same principal PA194C table.
Some additionally contain PA194CS, PA194CZ, or an eight-fold cached
search input. Treating every delivery as a new machine would therefore
misdescribe the current frontier.

The recognizer reads complete numeric binary tables in \file{.tm},
\file{.txt}, \file{.tsv}, and \file{.csv} files, including recursively
nested archives. It checks two rows per control, unique source/read
pairs, binary symbols, L/R moves, and one undefined halt target. The
filed collection replaces earlier archive wrappers, so occurrence counts
are not directly comparable to an incoming-only traversal. To check
coverage, every one of the 606 byte identities in the prior full census
was recovered and its action-graph identity recomputed. The new census
adds 14 byte forms.

\begin{center}
\begin{tabular}{@{}lr@{}}
\toprule
Quantity & Fresh result\\
\midrule
Research packets & 41\\
Filename-selected file occurrences inspected & 4,570\\
Recognized complete numeric-table occurrences & 2,685\\
Byte-distinct complete tables & 620\\
Root-free action-graph classes & 446\\
194-state occurrences / byte forms / graph classes & 95 / 15 / 8\\
\bottomrule
\end{tabular}
\end{center}

These are artifact counts, not numbers of correct consistency machines.
Two disconnected synthetic controls with six and nine states are excluded
from root-free classification; they are the only recognized sub-194
tables and are not proposed PA machines. The recognizer does not recover
arbitrary transition systems embedded only in program source or JSON.
The proposed frontier tables are present literally. The full member
ledger, frozen recognized tables, and comparison code are included.

\subsection{The eight existing 194-state graphs}

\begin{center}\small
\begin{tabularx}{\textwidth}{@{}lrrrX@{}}
\toprule
Design & Sites & Outcomes & Startup & Distinguishing construction\\
\midrule
PA194-67 & 3,673 & 4,780 & 58,154 & 67 readers, seven compatible folds.\\
PA194-67-AD & 3,673 & 4,780 & 58,122 & Physical exchange of A/D register roles.\\
PA194-68 & 3,673 & 4,780 & 58,154 & 68 readers, eight compatible folds.\\
PA194C & 3,649 & 4,749 & 57,922 & Cached MP conclusion and cached-head Shift.\\
PA194CS & 3,646 & 4,745 & 57,768 & Symmetric Shift removes a restoration loop.\\
PA194CZ & 3,649 & 4,750 & 58,168 & One zero-register jump at source site 26.\\
PA194-altC & 3,649 & 4,749 & 57,922 & Cached controller with the eight-fold decoder.\\
Rejected scanner & --- & --- & --- & Too-short scanner misreads a legal source path.\\
\bottomrule
\end{tabularx}
\end{center}
The old runtime and startup numbers in this table are those recorded by
the supplied constructions; this task freshly checked table identities,
source comparisons and novelty, and freshly verified the delivered new
machine. It did not rerun every historical physical proof.

The complete ordinal instruction graphs of PA194-67 and PA194-68 are
equal. Their seven and eight folds belong to different reader allocations;
they cannot be accumulated as fifteen independent savings. PA194C has
exactly the earlier PA195C source. PA194CS agrees with PA194C through
ordinal 3643 and replaces its last five-site Shift suffix by
\([\operatorname{inc}5;\operatorname{jump}0]\). Its extra reader is
balanced by the contextual identification \(41\to330\).

PA194-altC is easy to miss: it is a checked search input inside
\file{pa194c_conpa.zip}, rather than the packet's headline result. It
belongs in a novelty test nonetheless. The rejected scanner graph also
belongs in a byte/graph census, but not in a list of valid constructions.
Its source address at ordinal 1003 has seven consecutive zeroes, which
the shortened scanner mistakes for header padding. A local fold proof
alone does not prevent that global error.

\subsection{Larger families and transfer of ideas}

The U208 packets exploit a finite-history identification of controls
172 and 211 in a U209 ancestor. Many deliveries are alternate
presentations of the same graph. The relevant controls have already
disappeared in the PA196 development, so their saving cannot simply be
applied once more to PA194. James Chang's PA196 source has 3,812 primitive
sites; the later logical emitters and their layout searches account for
the transition to the current frontier.

The main 195-state families illustrate different routes to the same
count. Original PA195 uses a 191-token bundled diagonal emitter and
3,700 sites. PA195S separates the diagonal clause, using 190 tokens and
3,673 sites, with \(134+67-6\) controls. PA195C adds the cached controller,
reducing the source to 3,649 sites. PA195DC combines the bundled emitter
with that controller and has 3,676 sites. M195S8 achieves eight folds but
pays for an additional scanner control. These tradeoffs explain why a
source simplification or a larger fold count alone is insufficient.

NEG18 is another useful example. Its arithmetic identity shortens
encoded negation from 47 primitive sites to 18. In the separate-head
source, replacing eight occurrences reduces 3,673 sites to 3,441,
yet its published physical implementation has
197 states. Reader sharing, legal return paths and compatible local
folds dominate the final cost. The new searches in this packet included
mixed negation spellings on the newer grammar seeds; none supplied a
certified sub-194 table.

Appendix~\ref{app:ledger} gives an archive-by-archive comparison of all
41 packets. \file{comparison/archive_ledger.csv} also records exact
locators, hashes and evidence qualifications.

\section{The new construction}
\label{sec:construction}

\subsection{Registers, source instructions and the common seed}

The nine logical registers, numbered from zero, are
\[
 (N,T,W,A,D,H,X,Y,L).
\]
They are stored at physical unary positions
\((6,2,0,4,3,7,8,5,1)\), respectively. A source instruction is one of
\(\operatorname{inc}(r,u)\), \(\operatorname{dec}(r,z,p)\),
\(\operatorname{jump}(u)\), or halt. A decrement takes successor \(z\)
and leaves the registers unchanged when register \(r\) is zero;
otherwise it subtracts one and takes successor \(p\).

The clean source seed shared by PA194CS and PA194J is
\[
 c_*=(25,(2,0,0,0,0,0,0,0,0)).
\]
The source path of site 25 is \file{1000000}. The physical clean frame
has native control 84, head at \(-2\), header distance 48, the nine unary
registers above, and 26 spare blocks of value one. Both machines are
initialized from a literal blank tape; their different finite startup
traces are verified separately.

\subsection{Ten replacements, implemented by three rows}

\begin{center}\small
\begin{tabular}{@{}rrrrr@{}}
\toprule
Site & Zero register & Old jump / new zero target & New positive target & Masks\\
\midrule
3 & A (3) & 5 & 4 & 32\\
9 & A (3) & 11 & 10 & 16\\
26 & X (6) & 34 & 27 & 32\\
3526 & A (3) & 3528 & 3527 & 32\\
3532 & A (3) & 3535 & 3533 & 16\\
3544 & A (3) & 3546 & 3545 & 16\\
3550 & A (3) & 3553 & 3551 & 8\\
3553 & D (4) & 3637 & 3554 & 8\\
3627 & A (3) & 3629 & 3628 & 16\\
3633 & A (3) & 3637 & 3634 & 8\\
\bottomrule
\end{tabular}
\end{center}

``Masks'' counts the statuses occurring at that site in the finite
inductive abstraction defined below, not the number of possible natural
register vectors. All source paths are retained. Every instruction
outside these ten sites is identical, including its complete successor
list. The new source has 1,109 decrement sites, hence
\(3646+1109=4755\) primitive outcomes.

In native labels, the complete difference from PA194CS is
\begin{center}
\begin{tabular}{@{}rrll@{}}
\toprule
Control & Read & Parent action & PA194J action\\
\midrule
27 & 1 & \(0\,L\,40\) & \(1\,R\,21\)\\
103 & 1 & \(0\,L\,302\) & \(1\,R\,17\)\\
2066 & 1 & \(0\,L\,302\) & \(1\,R\,144\)\\
\bottomrule
\end{tabular}
\end{center}
Controls 21, 17 and 144 enter the respective decrement interfaces.
The last two changes provide reader pairs \((108,17)\) at control 103
and \((221,144)\) at control 2066. The first provides \((39,21)\) at
control 27. A pure reader preserves the scanned bit and moves right;
these pairs are useful as possible helper interfaces in subsequent
history-based folds. Their availability has not yet yielded a saved state.

The eight shared replacements arise from the source-tree definition
\file{node1 = (node0, breakTo 1)}, changed to
\file{node1 = (node0, dec 3)}. The scheduler definition changes
\file{node22 = (dec 8, breakTo 1)} to \file{(dec 8, dec 6)}.
The unique final AX \file{breakTo 1} becomes \file{dec 4}.
The builder regenerates the whole 3,646-site source from these tree edits;
it does not merely patch a supplied successor list.

\subsection{Why a reachable invariant is essential}

The replacement at site 9 has positive successor 10, a halt site. Thus
the two source programs are not interchangeable on arbitrary register
inputs. The proof must establish that the positive branch is unreachable
from the common seed. It does so for all ten replacements simultaneously.
In contrast, the \emph{physical simulation} theorem checks every branch
for arbitrary registers, including the positive branches. These two
different scopes fit together: the source invariant selects the actual
branches, while the physical theorem implements the complete new source.

\section{An inductive proof of exact source-run preservation}
\label{sec:invariant}

For a natural register vector \(r=(r_0,\ldots,r_8)\), define its positivity
mask
\[
 \alpha(r)=\sum_{j=0}^{8}2^j[\,r_j>0\,]\in\{0,\ldots,511\}.
\]
Write \(m\setminus j\) for clearing bit \(j\) and \(m\cup j\) for
setting it. Abstract states are pairs \((i,m)\) of a source ordinal and a
mask. For an instruction at site \(i\), define the abstract successors by
\[
\begin{array}{c|l}
\text{Instruction}&\text{Abstract successors}\\ \hline
\operatorname{jump}(u)&\{(u,m)\}\\
\operatorname{inc}(j,u)&\{(u,m\cup j)\}\\
\operatorname{dec}(j,z,p),\ m_j=0&\{(z,m)\}\\
\operatorname{dec}(j,z,p),\ m_j=1&\{(p,m),(p,m\setminus j)\}\\
\mathsf{halt}&\varnothing.
\end{array}
\]
The two positive-decrement possibilities reflect the distinction between
an input equal to one and an input greater than one. Discarding that
distinction is a deliberate overapproximation.

\begin{lemma}[Sound abstraction]
Every concrete source step \((i,r)\to(i',r')\) projects to an allowed
abstract step \((i,\alpha(r))\to(i',\alpha(r'))\).
\end{lemma}
\begin{proof}
Jumps preserve the vector. An increment makes its component positive.
A decrement of zero preserves the vector and takes the zero target.
If the component is positive, subtracting one either makes it zero or
leaves it positive, with every other bit unchanged. These cases are
exhaustive over the natural numbers.
\end{proof}

Let \(P\) be the parent PA194CS source and \(P'\) the PA194J source.
Compute the least abstract closure \(R\) of \((25,1)\) under \(P\).
The included builder, a separate finite checker, and an independent
review implementation reproduce the following finite facts:
\begin{enumerate}
\item \(|R|=263180\), and \((25,1)\in R\).
\item The abstract successors of every element of \(R\) lie in \(R\).
\item At every element of \(R\), the \(P\) and \(P'\) abstract successor
      sets agree exactly. Recomputing the least closure for \(P'\) gives
      the same \(R\).
\item At each changed site, the chosen register bit is zero in every
      mask in \(R\) at that site.
\end{enumerate}
The canonical sorted-pair certificate has SHA-256
\begin{quote}\small
\nolinkurl{9191fff6d156190b0eaa97f93dd5d10a23fbeac2d6505c0a269108e91c86f483}.
\end{quote}
The certificate is a finite object, not a bounded simulation of a source
run. It says nothing about a bound of 263,180 execution steps; the
invariant applies at every step, for arbitrarily large register values.

\begin{lemma}[Invariant and zero tests]
Every source configuration reachable from \(c_*\), in either program,
projects into \(R\). Consequently every inserted decrement takes its
zero branch and leaves every register unchanged.
\end{lemma}
\begin{proof}
The seed projects to \((25,1)\). Induct on the length of a concrete
source run, using sound abstraction and finite closure for the inductive
step. At a modified site, the corresponding bit in the projected mask
is zero, which is equivalent to the natural register being zero.
\end{proof}

\begin{theorem}[Exact source identity from the seed]
For every finite number of source steps, \(P\) and \(P'\), started at
\(c_*\), have identical source ordinals and identical nine-register
vectors, until and including any halt. In particular,
\[
 \PA\vdash\bigl(\Halt(P,c_*)\leftrightarrow\Halt(P',c_*)\bigr).
\]
\end{theorem}
\begin{proof}
Induct on steps while using the preceding invariant. At an unchanged
site the concrete instructions agree. At a changed site the parent
jumps to \(u\), while the child decrements a register equal to zero and
takes precisely \(u\), preserving the vector. The halt sites themselves
are unchanged. Thus the concrete states agree at the next step.

All definitions involved are primitive recursive. The finite closure
claims can be expressed using a fixed coded finite set and checked by
bounded computations in PA. The abstraction lemma is elementary
arithmetic, and the remaining argument is ordinary induction on run
length. The proof therefore establishes the displayed internal PA
statement, not only standard-model agreement.
\end{proof}

The source proof depends on the seed. It is neither an all-tape
equivalence between the literal machines nor an arbitrary-register
equivalence between the source programs. No such stronger statement is
needed for the blank-start theorem.

\section{Binding the source to the literal binary table}
\label{sec:physical}

\subsection{Complete symbolic frames}

A source frame represents a source path, zero padding, the header
delimiter, nine unary registers and a finite nonempty list of spare
unary blocks. The block for a value \(v\) is
\(B(v)=0\,1^{v+1}\). The checker also uses
\(C(v)=1\,0\,1^v\) during sweeps. It represents the complete tape,
including the zero exterior, rather than a finite observed window.
The base header distance is 48; its extra gap is an arbitrary natural.
The primary values, first spare value and remaining spare-list width are
independent symbolic quantities.

The source map is resolved twice, by an arithmetic path resolver and an
interval resolver. Both give exactly the generated instruction list.
This matters because equality of a flattened token word alone need not
preserve the targets of structured \file{breakTo} and \file{continueTo}
instructions.

\subsection{Universal physical contracts}

The independent checker derives four generic list-sweep lemmas from the
literal transition rows, for arbitrary values and arbitrary exterior
contexts. Finite-list induction extends the one-block calculations to
all finite spare lists. Literal steps, monotone unary-loop summaries,
these list lemmas, and exact word normalization then prove the source
contracts.

For each increment or jump there is one outcome; for each decrement,
there are separate zero and positive outcomes. A nonhalting outcome must
reach the precise target source frame with the required register effect
and spare-list update. A halt outcome must reach the terminal control.
Every duration is a positive affine expression in the symbolic
parameters. Consequently the simulation is productive: a source step
cannot be represented by a zero-length physical loop.

\begin{center}
\begin{tabular}{@{}lr@{}}
\toprule
Check & Result\\\midrule
Working states / transition rows & 194 / 388\\
Primitive source sites / outcomes & 3,646 / 4,755\\
Fresh symbolic rule applications & 1,103,582\\
Maximum symbolic steps per case & 403\\
Minimum actual-duration constant & 3\\
Generic list lemmas derived from literal rows & 4\\
Arithmetic-box oracles used & 0\\
Independent source resolvers in agreement & 2\\
Exact blank startup transitions & 60,062\\
\bottomrule
\end{tabular}
\end{center}
The complete physical replay returned \textbf{PASS}. The canonical contract
list has SHA-256
\begin{quote}\small
\nolinkurl{6c49226552b828b9d0325f358df908dc83995b0e08e851b38900db3c08124b9e}.
\end{quote}


The checker imports no allocator or abstract arithmetic-box oracle.
It is still ordinary executable mathematical evidence, not a proof
assistant kernel. The inherited symbolic engine, its normalization
identities and finite-list proof rules are exposed as readable source
under \file{construction/reference/}, with their provenance and licenses.

\subsection{Exact blank initialization}

A separate bytearray simulator executes precisely 60,062 literal
transitions from native control 74 on the all-zero tape. It obtains
native control 84, head \(-2\), and the frame of \(c_*\):
\[
 \text{path}=1000000,\quad\text{header distance}=48,\quad
 (N,T,W,A,D,H,X,Y,L)=(2,0,0,0,0,0,0,0,0).
\]
There are 26 spare blocks of value one and exactly 66 one-cells in the
whole tape. Every cell outside the explicit frame is zero. A set-based
simulator independently reproduces the same finite-support tape.
Thus initialization is a complete configuration equality, not a check of
only the register values or a small neighborhood of the head.

\begin{proposition}[Physical halting equivalence]
For PA194J, PA proves
\[
 H_\M\leftrightarrow\Halt(P',c_*).
\]
\end{proposition}
\begin{proof}
The finite initialization reaches the represented source seed without
halting. Concatenate the terminating positive physical simulations of
source instructions. A source halt gives a finite physical halting trace.
Conversely, suppose a physical halt occurs after \(t\) further transitions.
If no source halt occurred, the first \(t+1\) source instructions would
have terminating nonhalting simulations, each of positive duration.
Their concatenation extends beyond time \(t\), contradicting the alleged
halt. Intermediate tape configurations need not have a canonical source
decoding for this argument. The explicit finite traces, unary and
finite-list inductions, and induction on concatenated source steps are
all formalizable in PA.
\end{proof}

The same established interface for the supplied PA194CS source~\cite{cached} yields
the corresponding parent statement. Combining with
Section~\ref{sec:invariant} gives PA-provable blank halting equivalence
between PA194J and PA194CS.

\section{The inherited PA consistency calibration}
\label{sec:calibration}

It would be insufficient to infer PA-provability merely from a claim
that both machines do not halt in the standard model. The supplied
calibration instead provides syntactic transformations in both
directions. The relevant proof texts are retained in
\file{references/logical_calibration/}, with the cached-command correction
and the ten-entry selector adjustments. The inherited symmetric-guard
proof is retained separately as \file{proof_symmetric.tex}; it establishes
the precise cached command interface of PA194CS~\cite{cached}.

\subsection{Numerical formulas and commands}

Let
\[
 p(a,b)=\frac{(a+b)(a+b+1)}2+b,
 \quad I(A,B)=4p(A,B),\quad M(a,b)=4p(a,b)+1,
\]
\[
 U(a,B)=4p(B,a)+2,\qquad E(a,b)=4p(a,b)+3.
\]
The readings are implication, membership, universal quantification, and
equality. Code zero is truth; code one, \(M(0,0)\), is contradiction.
Negation is \(I(B,1)\). Formula-code bounds make the diagonal clause
\(\Gamma(B)=I(B,U(B,B))\) sound: variable \(x_B\) is absent from the
formula coded by \(B\). The current emitter pushes this clause
separately from the guarded pointed-separation gate.

In the cached command system, stores are \(\langle t,w,d,h\rangle\).
Shift halts exactly when \(h=1\); otherwise it maps the store to
\(\langle h,\operatorname{tl}(w),d,\operatorname{hd}(w)\rangle\).
For \(y=d+x\), modus ponens maps it to
\(\langle0,w,0,y\rangle\) if \(I(t,y)=h\), and to
\(\langle0,w,y,0\rangle\) otherwise. AX supplies ten pushed entries,
counting the saved top, and clears the counter and cache.

PA194CS implements the same command behavior through its symmetric
destructive equality routine. PA194J has exactly the same seeded
primitive execution as PA194CS, so it needs no changed command calculus,
selector, axiom schema or proof-word encoding.

The ten-entry command identities use
\[
\begin{aligned}
 \mathsf{Push}_{10}&=\mathsf{AX}_{0,0,0};S^9;\mathsf{MP}_0,\\
 \mathsf{Pop}_{10}&=\mathsf{AX}_{0,0,0};S^{11},\\
 \mathsf{Discard}_{10}&=\mathsf{MP}_1;\mathsf{AX}_{1,0,0};S^{10};\mathsf{MP}_0.
\end{aligned}
\]
The supplied command audit proves the full store effects, treats early
halts explicitly, and uses the total Cantor projection identities even
at \(p(0,0)=0\). In particular, the old nine-entry exponents cannot be
used unchanged. This interface has already been corrected in the parent
work and is preserved here.

\subsection{A source halt produces a PA refutation}

Write \(\operatorname{Bit}(s,z)\) for ``bit \(z\) of the finite set-code
\(s\) is one.'' The interpretation sends numerical membership to finite binary
membership \(\operatorname{Bit}(x_a,x_b)\), equality to arithmetic
equality, and \(U\) to universal quantification. PA proves that an
element of the finite set coded by \(s\) is smaller than \(s\), so it
refutes the interpretation of code one. Pointed separation is interpreted
using finite binary comprehension. For each arithmetic formula
\(\theta(z,\vec u)\), the induction proof of
\[
 \forall n\;\exists s<2^n\;\forall z<n\;
  (\operatorname{Bit}(s,z)\leftrightarrow\theta(z,\vec u))
\]
is obtained by a fixed proof-generating template. Substitution of a
formula's syntax into this template is primitive recursive, and PA
verifies its output as a proof. This is a statement about proof codes,
not a uniform truth predicate or a reflection principle.

Annotate the visible formula registers and canonical stack heads with
PA proofs of their interpretations. AX, MP and Shift preserve those
annotations. The zero-pair alias is handled by the truth annotation,
and parameter collisions in pointed separation use the corresponding
separate comprehension case. At a halt, the cache is code one, hence
the annotations yield a PA contradiction. The supplied argument is a
syntactic finite-trace transformation; it gives the PA-provable
implication from source halting to \(\neg\Con(\PA)\).

\subsection{A PA refutation produces a source halt}

The reverse translation recovers the required equality and membership
interfaces from the diagonal calculus. Reverse membership rewriting and
neighborhood equality use code freshness and diagonal vacuity;
the full proof is retained in \file{diagonal_calculus.tex} and
\file{original_pa195_calibration.tex}. Fixed proof templates turn the
bundled diagonal gate into the separate clause and gate, and back.
The correctly indexed ten-entry command compiler then translates a
refutation into a halting command word.

The strong form of this compiler is essential: its final word succeeds
from every incoming store, or an earlier halt has already occurred.
The scheduler retains stores between candidates, so a compiler requiring
a fresh empty store at each candidate would not suffice. The supplied
universal word begins with appropriate cleanup and handles every
preserved stack tail. The exhaustive scheduler eventually reaches that
word. All decoding, finite proof construction and simulation witnesses
are primitive recursive, with correctness proved in PA.

Consequently the parent source satisfies
\[
 \PA\vdash\bigl(\Halt(P,c_*)\leftrightarrow\neg\Con(\PA)\bigr).
\]
The new theorem follows by the two equivalences already proved:
\[
 H_\M\ \longleftrightarrow\ \Halt(P',c_*)\
 \longleftrightarrow\ \Halt(P,c_*)\
 \longleftrightarrow\ \neg\Con(\PA).
\]
This composition requires no assumption that PA is consistent. If PA is
consistent, actual nonhalting follows. No new end-to-end Lean/Rocq
elaboration of the composed theorem is asserted.

% Supplemental section; no preamble.  Uses the report's theorem environments
% and the existing bibliography keys pa194, cached, and cz.
\section{PA194D: a different logical emitter with 194 states}
\label{sec:gammad}

There is a second 194-state construction, denoted
\(\mathsf{PA194D}\), which changes the numerical axiom command itself.
The separate diagonal clause in the earlier machines is
\(\Gamma(A)=A\to U(A,A)\).  PA194D supplies
\(\Gamma(D)=D\to U(D,D)\) instead, while preserving the parameter
choices in all six other heads and in the guarded separation gate.
The resulting proof system is the same numerical calculus, but its
command interface and its physical reader allocation are different.
Its state count is
\[
 134\text{ fixed controls}+69\text{ raw readers}
       -9\text{ compatible folds}=194.
\]
The command controller is the cached controller with the scheduler's
known-zero decrement, as in the supplied PA194CZ construction
\cite{cached,cz}.  In addition, the final AX return uses a decrement of
the cleared register \(D\).  The argument below separates the logical
change from these two reachable source specializations.

The dense binary table has working controls \(0,\ldots,193\), blank
start control 0, and terminal halt label 194.  Its SHA-256 is
\begin{quote}\small
\nolinkurl{4226b7702de554df4b52d10239b2bd7ad79ca90bd8faf2bd256552fe0c7fe81b}.
\end{quote}
The native-label table starts at 74, halts at 254, and has SHA-256
\begin{quote}\small
\nolinkurl{bc8ab6c58badd137535938c9810e39d19152c49dd2eb26ef4ea07dbcf3157231}.
\end{quote}
The table uses the same two-way infinite binary tape convention as
PA194J. It has 388 defined transition rows. The authoritative file is
\file{logical_emitter/machine/pa194d.tm}; its complete contents are printed
in Appendix~\ref{app:table_d}.

\subsection{The exact logical change}

Use Cantor pairing \(p\), implication \(I\), membership \(M\),
universal quantification \(U\), and equality \(E\) with the numerical
definitions in Section~\ref{sec:calibration}.  In particular,
\(U(i,B)=4p(B,i)+2\), \(\neg B=I(B,1)\), and
\(\exists x_i B=\neg U(i,\neg B)\).  Write \([a;w]=p(a,w)\),
with the leftmost entry of a displayed list at its top.  The six retained
heads, in their original push order, are
\begin{align*}
 h_1(D,H)&=(\exists x_D U(D,H))\to H,\\
 h_2(D,A,H)&=U(D,U(A,H))\to U(A,U(D,H)),\\
 h_3(D,A,H)&=M(H,D)\to(E(H,A)\to M(A,D)),\\
 h_4(D,A,H)&=E(H,D)\to(M(A,H)\to M(A,D)),\\
 h_5(D,H)&=\exists x_D E(H,D),\\
 h_6(D,A,H)&=U(D,A)\to(U(D,A\to H)\to U(D,H)).
\end{align*}
The dominated propositional axiom is
\[
 \rho_8(D,A,H,u)=((H\to u)\to D)
                         \to((D\to H)\to(A\to H)),\qquad u\le D.
\]
Let \(G(a,b,c)\) be the retained pointed-separation formula from the
supplied separate-head calculus, and put
\[
 K(a,b,c)=
 \begin{cases}G(a,b,c),&c\ne a\text{ and }c\ne b,\\
 0,&\text{otherwise}.
 \end{cases}
\]
For precision, its construction is
\begin{align*}
 V(a,b,c)&=M(a,b)\to U(c,M(a,b)),\\
 P(a,b,c)&=c\to M(c,b),\\
 S_*(a,b,c)&=U\bigl(b,(M(a,b)\to P(a,b,c))
                                \wedge(M(c,b)\to c)\bigr),\\
 Q_*(a,b,c)&=U(c,S_*(a,b,c)\to\neg P(a,b,c)),\\
 G(a,b,c)&=\neg(V(a,b,c)\to Q_*(a,b,c)),
 \qquad B\wedge C=\neg(B\to\neg C).
\end{align*}

For an incoming store \(\langle t,w,d,h\rangle\), define
\(A=c\), \(D=d+x\), and \(H=h+y\).  The new command is exactly
\[
\begin{aligned}
 \mathsf{AX}^{D}_{c,x,y}\langle t,w,d,h\rangle
   &=\langle K(H,D,A),W',0,0\rangle,\\
 W'&=[\Gamma(D),h_6,h_5,h_4,h_3,h_2,h_1,\\[-2pt]
    &\hspace{3.5em}\rho_8(D,A,H,d),U(d,h),t;w].
\end{aligned}
\tag{D-AX}\label{eq:gammad-ax}
\]
Every head uses the displayed \(D,A,H\); only the argument of the
separate \(\Gamma\) has changed.  Thus the command still pushes ten
entries, counting the saved top, and still clears \(d,h\).
A failed gate guard does not remove the separately pushed
\(\Gamma(D)\).

This change is visible even on a zero incoming store.  For the command
parameters \((c,x,y)=(0,1,0)\), the former first stack entry is
\(\Gamma(0)=20\), whereas the new one is \(\Gamma(1)=832\).
The other emitted entries retain their former values.  In particular,
this is a change of the correlation between axiom parameters, not a
global exchange of physical register names.

Let \(\mathcal C_s\) be the supplied separate-head calculus: the six
families above, dominated \(\rho_8\), every \(\Gamma(B)\), each
guarded \(G(a,b,c)\), modus ponens, and theorem generalization
\(B\mapsto U(i,B)\).  ``Theorem'' means that there are no undischarged
hypotheses; the formula may have free variable labels.  The inherited
calibration gives primitive recursive proof translations, with
correctness provable in PA, establishing
\[
 \mathrm{PA}\vdash
 \left(\exists q\,\operatorname{Prf}_{\mathcal C_s}(q,1)
                   \ \leftrightarrow\ \neg\operatorname{Con}(\mathrm{PA})\right).
 \tag{D-Cal}\label{eq:gammad-cal}
\]
This is the separate-head calibration retained in
\file{references/logical_calibration/separate_head_calculus.tex},
together with its diagonal-calculus and PA195 dependencies
\cite{pa194}.  PA194D uses precisely this calculus.  It requires neither
a new logical axiom nor a new translation of PA proofs.

\subsection{All-store command identities and the changed selector}

For this subsection, use the ordinary cached-C source with the new
emitter and with its two original unconditional jumps.  Its numerical
commands have arbitrary stores \(\langle t,w,d,h\rangle\); scratch
registers obey the inherited clean command-boundary convention.  Write
\(S\) for cached Shift:
\[
 S\langle t,w,d,h\rangle=
 \begin{cases}
 \mathsf{halt},&h=1,\\
 \langle h,\operatorname{tl}(w),d,\operatorname{hd}(w)\rangle,&h\ne1.
 \end{cases}
\]
For \(y=d+x\), cached modus ponens is
\[
 \mathsf{MP}_x\langle t,w,d,h\rangle=
 \begin{cases}
 \langle0,w,0,y\rangle,&I(t,y)=h,\\
 \langle0,w,y,0\rangle,&I(t,y)\ne h.
 \end{cases}
\]
These are exactly the commands proved for the supplied cached
controller \cite{cached}.

Number the generated stack entries in (\ref{eq:gammad-ax}) by
\(q_1,\ldots,q_9\), so \(q_1=\Gamma(D)\),
\(q_8=\rho_8(D,A,H,d)\), and \(q_9=U(d,h)\).
The first eight entries and the gate have residue zero modulo four;
\(q_9\) has residue two.  None is code one.  Consequently the
ten-entry stack calculations in the supplied cached command audit apply
without any change in their exponents.  Explicitly, define
\begin{align*}
 \mathsf{Push}&=\mathsf{AX}^{D}_{0,0,0};S^9;\mathsf{MP}_0,\\
 \mathsf{Pop}&=\mathsf{AX}^{D}_{0,0,0};S^{11},\\
 \mathsf{Discard}&=\mathsf{MP}_1;\mathsf{AX}^{D}_{1,0,0};S^{10};\mathsf{MP}_0.
\end{align*}
Their exact contracts are
\begin{align*}
 \mathsf{Push}\langle t,w,d,h\rangle
   &=\langle0,[t;w],0,0\rangle,\\
 \mathsf{Pop}\langle t,[a;w],d,h\rangle
   &=\begin{cases}\mathsf{halt},&t=1,\\
                    \langle t,w,0,a\rangle,&t\ne1,
      \end{cases}\\
 \mathsf{Discard}\langle t,w,0,0\rangle
   &=\langle0,w,0,0\rangle.
\end{align*}
For example, \(\mathsf{AX}^{D};S^9\) leaves
\(\langle q_8,[t;w],0,q_9\rangle\), and
\(\mathsf{AX}^{D};S^{10}\) leaves
\(\langle q_9,w,0,t\rangle\).  The first ten Shifts test only zero
and generated entries; the eleventh tests the saved top.  Push's last
MP fails because an implication has residue zero and \(q_9\) has
residue two.  Discard's first MP fails since \(I(t,1)>0\); after
its AX and Shifts, the last MP fails since \(I(U(1,0),0)>0\).
These proofs use the total projection equations
\(\operatorname{hd}[a;w]=a\) and
\(\operatorname{tl}[a;w]=w\), including the alias \([0;0]=0\).
They do not presume that a zero push increases a stored stack length.

For \(1\le j\le7\), define a selector whose argument order is
\emph{\(D,H,A\)}:
\[
 \mathsf{select}_{j}(D,H,A)=
 \mathsf{Discard};\mathsf{AX}^{D}_{A,D,H};S^{j+1};
                                      \mathsf{Pop}^{9-j}.
\]
It sends every \(\langle t,w,0,0\rangle\) to
\(\langle q_j,w,0,0\rangle\).  The desired change to the proof
compiler is therefore exactly
\[
 \boxed{\quad
 \Gamma(B):\qquad\mathsf{select}_{1}(B,0,0)
       =\mathsf{Discard};\mathsf{AX}^{D}_{0,B,0};S^2;\mathsf{Pop}^{8}.
 \quad}
 \tag{D-Sel}\label{eq:gammad-selector}
\]
After Discard, \(d=h=0\), so this AX has \(D=B\).  The former
selector used \(\mathsf{select}_1(0,0,B)\).  Selection is safe even
when \(B=0\) and even when the unrelated gate guard fails.
The six retained heads use \(j=2,\ldots,7\), with their original
parameters and words.

For completeness, let \(\mathsf{Offset}_0\) be empty and let
\(\mathsf{Offset}_u=\mathsf{MP}_u\) for \(u>0\); from
\(\langle0,w,0,0\rangle\) it produces
\(\langle0,w,u,0\rangle\).  For \(u\le a\), the word
\[
 \mathsf{Discard};\mathsf{Offset}_u;
 \mathsf{AX}^{D}_{b,a-u,c};S^9;\mathsf{Pop}
\]
selects \(\rho_8(a,b,c,u)\).  The word
\(\mathsf{Discard};\mathsf{AX}^{D}_{c,b,a};\mathsf{Pop}^{10}\)
selects \(K(a,b,c)\).  Finally
\[
 \mathsf{Gen}_i=\mathsf{Push};\mathsf{Offset}_i;S;
                                      \mathsf{AX}^{D}_{i,0,0};S^{10}
\]
sends \(\langle a,w,0,0\rangle\) to
\(\langle U(i,a),w,0,0\rangle\), also when \(a=1\).
The Shift before AX tests zero and puts \(a\) in the cache; AX
clears that cache before a later Shift can test it.

\subsection{A strong proof compiler and both logical directions}

The cached-MP suffix is
\[
 \mathsf{MPSuffix}_b=S;S;\mathsf{MP}_b;S.
\]
From \(\langle0,[a,I(a,b),0;w],0,0\rangle\), it halts if
\(a=1\) or \(b=1\), and otherwise returns
\(\langle b,w,0,0\rangle\).  The second Shift tests \(a\);
successful MP places \(b\) in the cache; the final Shift tests \(b\)
and retrieves the dummy zero.  The dummy is necessary to preserve
the underlying \(w\).  This calculation also holds at the numerical
zero-stack alias.

Compile axiom leaves with the selectors just given, and compile a
generalization node by appending \(\mathsf{Gen}_i\) to its premise
word.  If an MP node has minor proof \(p_0\) of \(a\) and major
proof \(p_1\) of \(I(a,b)\), use
\[
 W_p=\mathsf{Discard};\mathsf{Push};
     W_{p_1};\mathsf{Push};W_{p_0};\mathsf{Push};\mathsf{MPSuffix}_b.
\]

\begin{lemma}[Uniform Gamma-D proof compiler]
For every finite \(\mathcal C_s\)-proof \(p\) and every \(t,w\),
the word \(W_p\) from \(\langle t,w,0,0\rangle\) either halts
or returns \(\langle\operatorname{concl}(p),w,0,0\rangle\).
If \(p\) proves code one, then
\[
  \mathsf{AX}^{D}_{0,0,0};W_p;\mathsf{Pop}
\]
halts from every incoming four-component store.
\end{lemma}
\begin{proof}
Induct on the proof tree.  The selectors prove the axiom case, including
the changed selector (\ref{eq:gammad-selector}).  The generalization
contract proves that case.  In the MP case, Discard and Push introduce
a dummy zero; if neither recursive premise word halts, their contracts
and Push produce precisely the suffix's three-entry precondition.
The suffix gives either an early halt or the required conclusion with
the original tail.  An earlier halt remains a halt when a word is
extended.  For the second claim, the initial AX clears both caches,
so the first claim applies to its arbitrary resulting top and stack.
If no earlier halt occurs, the last Pop tests the derived top one
before inspecting the arbitrary underlying stack.
\end{proof}

This all-store clause matters because the inherited exhaustive
command-word scheduler retains its stores between trials.  A refuting
proof yields a finite universally halting word, so the scheduler either
halts earlier or halts when it reaches this word.  The inherited
word encoding and scheduler proof are unchanged.

Conversely, annotate \(t\), \(h\), and each decoded stack entry
with a \(\mathcal C_s\)-derivation.  Code zero has the fixed proof
\(\rho_8(0,0,0,0)=0\), which supplies all zero padding.
AX preserves these annotations: the six heads and a successful gate
are axioms; a failed gate is zero; \(\Gamma(D)\) is an unconditional
instance of the existing diagonal family; \(d\le d+x\) permits
the propositional head; and \(U(d,h)\) follows by generalization
from the old cache.  A successful MP derives \(d+x\) from the
annotations for \(t\) and \(h=I(t,d+x)\); a failed MP returns
zero.  Shift moves derived entries, and a halting Shift exposes a
proof of one.  No derivability condition is required of the offset
\(d\).  The actual seeded source begins with the required zero
annotations.  This proves that its halting yields a refutation.

All selectors and proof-word constructors are primitive recursive.
The arguments are inductions on proof trees, finite lists, source loop
counters, and traces.  Together with the unchanged scheduler they prove
in PA that the ordinary Gamma-D cached-C source halts exactly when
\(\mathcal C_s\) has a refutation.  The proof translations in
(\ref{eq:gammad-cal}) then give its equivalence to
\(\neg\operatorname{Con}(\mathrm{PA})\).  In particular, the
soundness direction does not appeal to the standard truth of arithmetic
or to reflection for PA.

\subsection{Exact source binding and the two zero tests}

Let \(P_D\) be that ordinary cached-C Gamma-D source and let
\(P_D^Z\) be the source implemented by PA194D.  Both have 3,649
primitive sites, numbered from zero.  Relative to the old 190-token
Gamma-A emitter, the copy macros at token indices 121, 122, and 123
change from \texttt{copyA} to \texttt{copyD}.  Their source spans
are respectively 2258--2283, 2284--2309, and 2310--2335.
Each span has 26 primitive instructions.  The six altered register
arguments, and all their exact successors, are
\begin{center}\small
\begin{tabular}{@{}rll@{}}
\toprule
Site & Old instruction & Gamma-D instruction\\
\midrule
2277 & \(\operatorname{dec}(3,2281,2278)\) & \(\operatorname{dec}(4,2281,2278)\)\\
2282 & \(\operatorname{inc}(3,2283)\) & \(\operatorname{inc}(4,2283)\)\\
2303 & \(\operatorname{dec}(3,2307,2304)\) & \(\operatorname{dec}(4,2307,2304)\)\\
2308 & \(\operatorname{inc}(3,2309)\) & \(\operatorname{inc}(4,2309)\)\\
2329 & \(\operatorname{dec}(3,2333,2330)\) & \(\operatorname{dec}(4,2333,2330)\)\\
2334 & \(\operatorname{inc}(3,2335)\) & \(\operatorname{inc}(4,2335)\)\\
\bottomrule
\end{tabular}
\end{center}
Here registers 3 and 4 are logical \(A\) and \(D\), respectively.
Every other instruction of \(P_D\), including every successor,
agrees with the old cached-C source.  A separate free-expression
calculation checks the entire stack and gate of
(\ref{eq:gammad-ax}) for symbolic incoming parameters.

The physical search reassociates the emitter tree.  The winning
grammar's complete resolved instruction graph equals that of an
independently balanced tree over the same macro word.  It also equals
the graph obtained by changing only these macros in the supplied
accepted association.  Thus the binding checks structured jump
destinations, not only equality of the flattened token sequence.

To pass from \(P_D\) to \(P_D^Z\), change precisely
\[
\begin{array}{c|cc}
\text{site}&P_D&P_D^Z\\ \hline
26&\operatorname{jump}(34)&\operatorname{dec}(6,34,27)\\
3553&\operatorname{jump}(3637)&\operatorname{dec}(4,3637,3554).
\end{array}
\]
The first tests scratch \(X\); the second is the final AX return
and tests \(D\).  This adds two branch outcomes: \(P_D\) has
4,749, and \(P_D^Z\) has 4,751.

Use the positivity abstraction of Section~\ref{sec:invariant}, now
computed for \(P_D\) from the same clean seed
\[
 c_*=(25,(2,0,0,0,0,0,0,0,0)).
\]
Its least abstract closure \(R_D\) contains 263,244 location/mask
pairs.  A separate replay verifies that the seed belongs to it,
that it is closed, that the two source programs have exactly the same
abstract successors on it, and that the tested bits are zero at the
two changed sites.  There are 32 masks at site 26, all with bit 6
clear, and eight masks at site 3553, all with bit 4 clear.
The sorted-pair certificate has SHA-256
\begin{quote}\small
\nolinkurl{57c58dcae511fe13a6c200e7cfe4f4322370aec72b4a2fe5ba39b906fe62dbf8}.
\end{quote}
This is a fresh invariant for the changed formula emitter; it is not
an assumption that the Gamma-A invariant still applies.

The abstraction lemma and induction on execution length therefore
prove exact equality of the \(P_D\) and \(P_D^Z\) source runs
from \(c_*\), until and including any halt.  Each inserted decrement
sees zero, preserves every register, and takes the former jump target.
PA proves the finite closure facts and the two inductions, hence
\[
 \mathrm{PA}\vdash
 \bigl(\operatorname{Halt}(P_D,c_*)
          \leftrightarrow\operatorname{Halt}(P_D^Z,c_*)\bigr).
\]
This source equivalence is claimed from the specified seed.  The
arbitrary-store command calculations above concern \(P_D\);
they are not being inferred for arbitrary inputs to the specialized
source \(P_D^Z\).

\subsection{Literal realization and the PA theorem}

The nine folds are performed simultaneously on the 203-control raw
allocation. In native labels their replacement rows are:
\begin{center}\small
\begin{tabular}{@{}rrll@{}}
\toprule
Removed & Retained & New read-0 action & New read-1 action\\
\midrule
2004 & 2063 & $0\,L\,2049$ & $1\,R\,90$\\
2015 & 2023 & $0\,L\,2061$ & $1\,R\,2012$\\
2018 & 2012 & $0\,L\,2055$ & $1\,R\,90$\\
2022 & 2033 & $0\,L\,2002$ & $1\,R\,61$\\
2025 & 2046 & $0\,L\,2034$ & $1\,R\,2029$\\
2028 & 2039 & $0\,L\,2014$ & $1\,L\,2003$\\
2037 & 2059 & $0\,L\,2044$ & $1\,R\,2038$\\
2053 & 2011 & $0\,R\,98$ & $1\,L\,2019$\\
2064 & 2032 & $0\,L\,2017$ & $1\,L\,2041$\\
\bottomrule
\end{tabular}
\end{center}
Every incoming row establishes the required predecessor bit. None of the
folded controls is the initial control. The 36 local cases are checked
on the final simultaneous quotient: 14 take one step, 19 take three,
and three take five. The traces preserve the exterior tape and the full
required local effect; the supplied universal source contracts also
verify the resulting whole interpreter independently.

The literal table has an independently checked physical simulation
contract for each of the 4,751 primitive outcomes of \(P_D^Z\), including the positive
branches that the seeded source invariant excludes.  These contracts
range over arbitrary natural primary values and header gap and every
finite nonempty spare list.  The independent checker uses 1,091,154
symbolic steps, with a maximum of 385 for one case; every duration has
nonnegative coefficients and constant at least three.  It derives four
list lemmas from the literal rows and invokes no arithmetic-box oracle.
The nine compatible reader folds have 36 local all-tape cases.
The exact blank initialization takes 58,168
literal transitions to the complete frame for \(c_*\), in native
control 84 with head at \(-2\).  The frame has source path 1000000,
header distance 48, primary vector \((2,0,0,0,0,0,0,0,0)\), and
26 spare blocks of value one.

The source simulations terminate in positive time.  Consequently a
source halt produces a physical halt, and a physical halt cannot
precede a nonhalting concatenation of more positive source simulations
than the alleged halting time.  As in Section~\ref{sec:physical}, this
argument requires no inverse decoding of arbitrary intermediate tapes.
The finite startup and local certificates, and the unary, list and
trace inductions, are formalizable in PA.

\begin{theorem}[Gamma-D consistency machine]
Let \(H_D\) be the arithmetic sentence asserting a finite halting
trace of the specified PA194D table from its blank start.  Then
\[
 \mathrm{PA}\vdash
    \bigl(H_D\leftrightarrow\neg\operatorname{Con}(\mathrm{PA})\bigr),
 \qquad
 \mathrm{PA}\vdash
    \bigl(\neg H_D\leftrightarrow\operatorname{Con}(\mathrm{PA})\bigr).
\]
\end{theorem}
\begin{proof}
Compose, inside PA, the physical simulation, the seeded zero-test
equivalence, the soundness and strong proof compiler, and the supplied
proof calibration:
\[
\begin{split}
 H_D
 &\longleftrightarrow\operatorname{Halt}(P_D^Z,c_*)
 \longleftrightarrow\operatorname{Halt}(P_D,c_*)\\
 &\longleftrightarrow
       \exists q\,\operatorname{Prf}_{\mathcal C_s}(q,1)
 \longleftrightarrow\neg\operatorname{Con}(\mathrm{PA}).
\end{split}
\]
Classical logic in PA gives the nonhalting formulation.
\end{proof}

The theorem is an ordinary proof with explicit finite and symbolic
certificates, using the retained ordinary PA calibration.  It is not
a new proof-assistant kernel export or an iPRA equivalence theorem.
It makes no consistency assumption in establishing the PA-provable
biconditional.  PA194D supplies a second construction at the requested
194-state fallback threshold; it does not lower that threshold.


\section{Novelty within the pinned corpus}
\label{sec:novelty}

For a deterministic binary action graph with a chosen root, breadth-first
discovery in read-bit order produces a canonical numbering. To forget
the root, the checker takes the lexicographically least serialization
over all controls from which the entire working graph is reachable.
The unique terminal control is marked separately. Two such graphs are
isomorphic exactly when these canonical serializations agree.

For the 194-state comparison, the calculation is repeated after global
L/R reflection, global interchange of tape symbols, and their combination.
Even allowing a change of designated start and the symbol-exchanged blank
convention, PA194J matches none of the eight existing 194-state classes.
Graphs of other cardinalities cannot be isomorphic to it. The four
variants choosing A, D, H or X at the final zero-test location are also
mutually distinct under the same tests; the delivered principal table
uses D.

The finite comparison establishes novelty relative to the retrieved
corpus. It is not a claim of worldwide priority. Its semantic significance
comes from the independently proved source changes and actual startup
use of changed rows, not from a hash inequality alone.

The independently audited PA194D graph and its CSZ sibling also differ
from all eight old graphs, from each other, and from all four new J-family
graphs under these same transformations. A separate free-term audit of
the entire emitter rules out every global permutation of \(A,D,H\):
exactly one of its ten pushed entries changes, while the other nine,
parameter reads and gate agree. The concrete head values 20 and 832 in
Section~\ref{sec:gammad} further expose that logical difference.

\section{Searches below 194 and their limits}

The following new searches jointly mutate source associations and evaluate
compatible finite-history folds. Closing-leaf searches additionally vary the
final jump between ordinary return and candidate zero-register decrements.
The score is the effective reader count plus scanner/address penalties;
it is a filter, not a physical correctness certificate. A score below 60
would be a candidate for fewer than 194 states with the 134-control fixed part.
Every such candidate would still require the complete source and tape proofs.

\begin{center}\small
\begin{tabular}{@{}lrrr@{}}
\toprule
Search & Mutations & Best score & Seconds\\
\midrule
CSZ & 11,263,380 & 60 & 900\\
CZ & 9,991,239 & 60 & 900\\
CSZ103A & 2,940,290 & 60 & 450\\
CSZ103X & 3,442,575 & 60 & 450\\
closing103A & 4,489,001 & 60 & 450\\
negation\_search & 3,062,902 & 60 & 450\\
h1D & 890,683 & 67 & 180\\
h1D\_adaptive & 1,128,819 & 64 & 240\\
separateD & 3,608,903 & 62 & 450\\
separateD\_adaptive & 1,927,245 & 60 & 240\\
separateH & 3,447,089 & 63 & 450\\
separateH\_adaptive & 4,398,267 & 63 & 240\\
warm\_D\_CSZ\_A & 9,288,577 & 60 & 480\\
warm\_D\_CSZ\_X & 9,325,004 & 60 & 480\\
\bottomrule
\end{tabular}
\end{center}
Logs, frozen input pools, random seeds and C++ source are retained under
\file{research/search/}. Wall-time stopping means mutation counts vary with
the host. No run produced a score below 60. The adaptive separate-D run
reached 69 readers and nine simultaneous folds; this candidate became
PA194D only after the independent proofs reported in Section~\ref{sec:gammad}.
Failed searches establish no lower bound and do not exhaust all machines,
logical calculi, controller layouts, or source associations.


The new work also explored changes to the logical emitter. Three
191-token nine-entry variants relocate a diagonal antecedent into a
quantifier head; two 190-token ten-entry variants emit \(\Gamma(D)\)
or \(\Gamma(H)\) instead of \(\Gamma(A)\). Exact free-algebra word
checks, primitive source graphs, proof arguments and discovery seeds are
under \file{research/logical_variants/}. The \(\Gamma(D)\) search reached
the second 194-state construction given above. The other four variants
remain logical source proposals and search inputs, without a claimed
194-state physical realization. PA194J uses the unchanged parent
numerical calculus; PA194D supplies the new selector calibration.

\subsection{Two concrete failures worth preserving}

Deleting scanner control 34 and redirecting it to 75 fails 113 of the
4,780 universal contracts for the supplied primary PA194 source. The failures include fixed
controller paths, so shortening AX addresses alone would not justify
the deletion. A finite blank trace that has not yet encountered a bad
address is insufficient evidence.

The ZFC zero-spare idea is also not a plug-in change. Altering native
row \(198/0\) from \(1L40\) to \(0L40\) in PA194CS fails its bootstrap:
it visits the source dispatcher only at times 38 and 66 before entering
an unbounded rightward one-writing drift. The ZFC construction had a
different scheduler/bootstrap proof and cannot supply the missing PA
initialization. This experiment is excluded from the delivered table.

\subsection{What transfers from ZFC and concise iPRA}

Useful ZFC methods~\cite{zfc} already present in the PA frontier include contextual
reader folds, source reassociation with exact control-flow binding,
cached conclusions, and exploiting destructive comparison symmetry.
The new construction extends the global zero-information method beyond
the single scheduler jump found in PA194CZ~\cite{cz}. Zero-spare and command-credit
ideas remain possible research directions only after a new bootstrap
and command proof are supplied.

ConciseIPRA~\cite{ipra} offers small proof presentations, primitive-recursive
syntax, and disciplined separation of logical and executable certificates.
Those ideas help organize a verifier, but replacing the checked theory
by iPRA would change the consistency statement. A smaller iPRA checker
does not directly give a smaller machine for \(\Con(\PA)\). The present
construction therefore preserves the PA-calibrated calculus.

\section{Reproduction and evidence boundaries}

Python 3.10 or later and its standard library suffice for the accepted
construction and checks. From the archive root, run
\begin{quote}\ttfamily
python3 -B verify.py
\end{quote}
Assertions must be enabled; the scripts reject optimization mode.
The full check rebuilds both literal tables and sources, recomputes the
finite invariants, runs both independent physical proofs, replays the whole
blank initializations, and checks corpus novelty and the changed logical
emitter. It works in temporary copies to preserve the frozen packet.
The optional \file{--quick} mode omits the universal physical replays;
it retains the other checks. Optional C++ searches
are not required. The saved receipts give independently reviewable
results but do not replace execution of the verifier.

To rebuild this article, run \file{build_report.sh}; it invokes
\file{pdflatex} twice with failure-on-error flags. The exact table,
source files, supporting proof texts and finite receipts are included.
\file{SHA256SUMS} binds the delivery. The attribution and license files
identify the inherited James Chang, ProveIt and symbolic-engine material.

The claims should be read at their actual levels:
\begin{itemize}
\item Fresh finite facts: state count, table hashes, source maps, graph
      distinctions, mask closure, and complete startup equality.
\item Fresh symbolic mathematics: arbitrary-register physical contracts,
      four list lemmas and positive simulation durations.
\item New ordinary theorem: exact seeded source preservation and the
      resulting PA-provable halting equivalence with the parent.
\item Inherited ordinary theorem: the numerical calculus, strong command
      compiler and scheduler are calibrated to PA consistency.
\item Unclaimed: a smaller machine, a general lower bound, a new
      end-to-end proof-assistant theorem, or an iPRA proof of the full
      consistency equivalence.
\end{itemize}

\section{Further research questions}

\begin{enumerate}
\item Can the new helper pairs \((39,21)\), \((108,17)\), and
      \((221,144)\) support a reader allocation of effective size 59?
      With the current 134-control fixed part, that would give 193 states.
      A search must include the source invariant and all physical branch
      contracts, not just the new pair identities.
\item Can one combine two different known-zero registers in a larger
      decoder gadget to replace three readers by two? The four certified
      final-register variants provide different interfaces for such a
      search. Longer detours would need a positive local simulation proof.
\item Can the fixed controller be reassociated to shorten \emph{all}
      return addresses, including startup and the equality guard, enough
      to remove one scanner control? AX-only address constraints miss
      the current fixed-path obstruction.
\item Can a smaller finite invariant or a structured loop invariant
      replace the 263,180-pair certificate? The current set is easy to
      recompute but is not optimized as a mathematical or Lean artifact.
\item Can the successful \(\Gamma(D)\) allocation of 69 readers and nine
      compatible folds be combined with the new PA194J helper pairs to
      obtain a tenth fold? Can the other four logical emitter variants
      gain comparable sharing? Their equiconsistency arguments alone do
      not determine the physical-state cost.
\item Can the new invariant theorem and the literal tape contracts be
      connected to the repository's existing Lean source and arithmetic
      endpoints? A first useful target is the seeded source equivalence,
      whose proof needs only finite closure and induction; the full iPRA
      calibration remains a separate task.
\end{enumerate}

\begin{thebibliography}{9}
\bibitem{repo} Vladimir Reshetnikov, \emph{ProveIt}, pinned revision
\nolinkurl{17f92105e2bc4317feae6c219cb2143437b482ff}.
\url{https://github.com/VladimirReshetnikov/ProveIt/tree/17f92105e2bc4317feae6c219cb2143437b482ff/Computability/BusyBeaver/Research/PA}.
\bibitem{pa194} Supplied PA194 research packets, 13 November 2026,
filed under \file{Computability/BusyBeaver/Research/PA/docs/pa194_conpa_2026-11-13}
and its \file{_2} companion. Exact source and table identities are in
the supplied comparison ledger.
\bibitem{cached} Supplied \emph{Two new 194-state PA consistency machines},
\file{docs/incoming/pa194_new_designs_2026-11-13.zip}, especially
\file{secondary/pa194CS} and \file{construction/logical_reference}.
\bibitem{cz} Supplied PA194C/PA194CZ packet,
\file{docs/incoming/pa194c_conpa_2026-11-13 (1).zip}, especially
\file{construction/build194cz.py}.
\bibitem{chang} James Chang, PA196 and PA209 developments, imported
snapshots under \file{PA/jamesyc-upstream}; original attribution and
licenses are preserved in the supplied sources.
\bibitem{zfc} ProveIt ZFC consistency-machine research, especially
\file{candidate-k215/alternatives/zero-credit} and its related source
and physical preservation arguments.
\bibitem{ipra} ProveIt, \emph{Concise presentations of intuitionistic
primitive recursive arithmetic},
\url{https://github.com/VladimirReshetnikov/ProveIt/tree/17f92105e2bc4317feae6c219cb2143437b482ff/Logic/PrimitiveRecursiveArithmetic/ConciseIPRA}.
\end{thebibliography}

\clearpage\appendix
\section{Complete dense transition table}
\label{app:table}
Each entry \(w\,D\,q'\) means: write \(w\), move in direction \(D\),
and enter \(q'\). The initial state is 0 and terminal halt label is 194.
The paired columns contain all 194 working states, with both read symbols.
\small
\begin{longtable}{@{}rll@{\hspace{1.3cm}}rll@{}}
\toprule
State & Read 0 & Read 1 & State & Read 0 & Read 1\\
\midrule\endfirsthead
\toprule
State & Read 0 & Read 1 & State & Read 0 & Read 1\\
\midrule\endhead
\midrule\multicolumn{6}{r}{Continued on the next page}\\\endfoot
\bottomrule\endlastfoot
0 & $1\,R\,1$ & $1\,R\,2$ & 97 & $0\,L\,112$ & $1\,R\,113$\\
1 & $1\,R\,3$ & $1\,R\,1$ & 98 & $0\,R\,114$ & $0\,L\,14$\\
2 & $0\,R\,0$ & $1\,R\,2$ & 99 & $0\,R\,115$ & $1\,R\,91$\\
3 & $0\,R\,4$ & $0\,R\,1$ & 100 & $0\,R\,115$ & $1\,R\,106$\\
4 & $1\,R\,5$ & $1\,L\,6$ & 101 & $0\,R\,116$ & $0\,L\,14$\\
5 & $1\,L\,7$ & $1\,L\,8$ & 102 & $0\,R\,68$ & $1\,L\,66$\\
6 & $0\,L\,7$ & $1\,L\,6$ & 103 & $0\,R\,117$ & $1\,R\,118$\\
7 & $0\,L\,9$ & $1\,L\,6$ & 104 & $0\,R\,89$ & $1\,R\,104$\\
8 & $0\,L\,10$ & $0\,L\,11$ & 105 & $0\,R\,14$ & $1\,R\,105$\\
9 & $1\,R\,12$ & $0\,L\,9$ & 106 & $0\,L\,112$ & $1\,L\,33$\\
10 & $0\,L\,9$ & $0\,L\,10$ & 107 & $0\,R\,67$ & $1\,R\,119$\\
11 & $1\,L\,8$ & $1\,L\,11$ & 108 & $0\,R\,63$ & $1\,R\,108$\\
12 & $0\,R\,12$ & $0\,R\,13$ & 109 & $0\,R\,120$ & $1\,L\,121$\\
13 & $0\,L\,14$ & $0\,R\,13$ & 110 & $0\,R\,122$ & $1\,R\,123$\\
14 & $1\,L\,15$ & $0\,L\,15$ & 111 & $0\,R\,46$ & $0\,L\,14$\\
15 & $0\,L\,15$ & $1\,L\,16$ & 112 & $0\,R\,86$ & $1\,R\,96$\\
16 & $0\,L\,17$ & $1\,L\,16$ & 113 & $0\,R\,112$ & $1\,L\,91$\\
17 & $0\,L\,18$ & $1\,L\,16$ & 114 & $0\,R\,124$ & $1\,R\,124$\\
18 & $0\,L\,19$ & $1\,L\,16$ & 115 & $0\,L\,60$ & $0\,R\,83$\\
19 & $0\,L\,20$ & $1\,L\,16$ & 116 & $0\,R\,62$ & $0\,R\,89$\\
20 & $0\,L\,21$ & $1\,L\,16$ & 117 & $0\,L\,81$ & $1\,R\,125$\\
21 & $0\,L\,22$ & $1\,L\,16$ & 118 & $0\,R\,126$ & $1\,R\,127$\\
22 & $0\,L\,23$ & $1\,L\,16$ & 119 & $0\,R\,47$ & $1\,R\,85$\\
23 & $0\,R\,23$ & $1\,R\,24$ & 120 & $0\,R\,55$ & $0\,R\,108$\\
24 & $0\,R\,25$ & $1\,R\,26$ & 121 & $0\,R\,35$ & $1\,R\,33$\\
25 & $0\,R\,27$ & $1\,R\,28$ & 122 & $0\,R\,95$ & $1\,R\,111$\\
26 & $0\,R\,29$ & $1\,R\,30$ & 123 & $0\,R\,75$ & $1\,R\,74$\\
27 & $0\,R\,31$ & $1\,R\,194$ & 124 & $0\,R\,46$ & $1\,R\,46$\\
28 & $0\,L\,32$ & $1\,R\,33$ & 125 & $0\,R\,128$ & $1\,R\,129$\\
29 & $0\,R\,34$ & $1\,R\,35$ & 126 & $0\,R\,32$ & $1\,R\,64$\\
30 & $0\,L\,33$ & $1\,L\,36$ & 127 & $0\,R\,31$ & $1\,R\,109$\\
31 & $0\,R\,37$ & $1\,R\,38$ & 128 & $0\,R\,130$ & $1\,R\,131$\\
32 & $0\,R\,39$ & $1\,R\,40$ & 129 & $0\,R\,132$ & $1\,R\,133$\\
33 & $0\,R\,41$ & $1\,R\,42$ & 130 & $0\,R\,134$ & $1\,R\,135$\\
34 & $0\,R\,43$ & $1\,R\,44$ & 131 & $0\,R\,136$ & $1\,R\,137$\\
35 & $0\,R\,45$ & $1\,R\,36$ & 132 & $0\,R\,138$ & $1\,R\,139$\\
36 & $0\,R\,46$ & $1\,R\,47$ & 133 & $0\,R\,140$ & $1\,R\,141$\\
37 & $0\,R\,48$ & $1\,R\,41$ & 134 & $0\,R\,93$ & $1\,R\,121$\\
38 & $0\,R\,49$ & $1\,R\,50$ & 135 & $0\,R\,142$ & $1\,R\,124$\\
39 & $0\,R\,51$ & $1\,R\,52$ & 136 & $0\,R\,143$ & $1\,R\,144$\\
40 & $0\,R\,53$ & $1\,R\,54$ & 137 & $0\,R\,145$ & $1\,R\,146$\\
41 & $0\,R\,55$ & $0\,R\,56$ & 138 & $0\,R\,134$ & $1\,R\,147$\\
42 & $0\,R\,57$ & $0\,L\,14$ & 139 & $0\,R\,46$ & $1\,R\,148$\\
43 & $0\,R\,58$ & $1\,R\,59$ & 140 & $0\,R\,149$ & $1\,L\,131$\\
44 & $0\,R\,60$ & $1\,R\,49$ & 141 & $0\,R\,150$ & $1\,R\,151$\\
45 & $0\,R\,55$ & $0\,R\,61$ & 142 & $0\,R\,72$ & $1\,R\,40$\\
46 & $0\,R\,62$ & $0\,R\,63$ & 143 & $0\,R\,28$ & $1\,R\,152$\\
47 & $0\,L\,64$ & $1\,L\,65$ & 144 & $0\,L\,153$ & $1\,L\,154$\\
48 & $0\,R\,45$ & $1\,R\,50$ & 145 & $0\,L\,129$ & $1\,R\,155$\\
49 & $0\,L\,66$ & $0\,L\,15$ & 146 & $0\,R\,156$ & $1\,R\,157$\\
50 & $0\,R\,67$ & $1\,R\,68$ & 147 & $0\,R\,149$ & $1\,R\,158$\\
51 & $0\,R\,69$ & $1\,R\,70$ & 148 & $0\,R\,159$ & $1\,R\,46$\\
52 & $0\,R\,71$ & $1\,R\,72$ & 149 & $0\,R\,93$ & $1\,R\,160$\\
53 & $0\,R\,73$ & $1\,R\,74$ & 150 & $0\,R\,161$ & $1\,R\,162$\\
54 & $0\,R\,75$ & $1\,R\,76$ & 151 & $0\,R\,163$ & $1\,R\,120$\\
55 & $0\,L\,77$ & $0\,R\,1$ & 152 & $0\,R\,164$ & $1\,R\,165$\\
56 & $0\,R\,61$ & $1\,R\,56$ & 153 & $0\,R\,166$ & $1\,R\,157$\\
57 & $0\,L\,69$ & $1\,L\,35$ & 154 & $0\,R\,167$ & $1\,R\,156$\\
58 & $0\,R\,78$ & $1\,R\,79$ & 155 & $0\,R\,135$ & $1\,R\,130$\\
59 & $0\,R\,51$ & $1\,R\,80$ & 156 & $0\,R\,147$ & $1\,R\,40$\\
60 & $0\,L\,81$ & $1\,R\,82$ & 157 & $0\,R\,134$ & $1\,R\,168$\\
61 & $0\,R\,83$ & $1\,R\,61$ & 158 & $0\,R\,40$ & $1\,L\,169$\\
62 & $0\,L\,77$ & $0\,R\,2$ & 159 & $0\,L\,170$ & $1\,R\,171$\\
63 & $0\,R\,84$ & $1\,R\,63$ & 160 & $0\,R\,102$ & $1\,R\,80$\\
64 & $0\,R\,30$ & $1\,R\,85$ & 161 & $0\,R\,172$ & $1\,R\,170$\\
65 & $0\,R\,86$ & $1\,R\,87$ & 162 & $0\,R\,173$ & $1\,R\,174$\\
66 & $0\,R\,35$ & $1\,R\,88$ & 163 & $0\,R\,175$ & $1\,R\,176$\\
67 & $0\,R\,41$ & $0\,L\,13$ & 164 & $0\,R\,177$ & $1\,R\,178$\\
68 & $0\,R\,55$ & $0\,R\,89$ & 165 & $0\,L\,179$ & $1\,L\,180$\\
69 & $0\,R\,90$ & $1\,R\,65$ & 166 & $0\,R\,149$ & $1\,R\,140$\\
70 & $0\,R\,47$ & $1\,R\,91$ & 167 & $0\,R\,165$ & $1\,R\,181$\\
71 & $0\,R\,41$ & $1\,R\,92$ & 168 & $0\,L\,130$ & $1\,R\,46$\\
72 & $0\,R\,93$ & $1\,R\,94$ & 169 & $0\,R\,182$ & $1\,R\,135$\\
73 & $0\,R\,95$ & $1\,L\,96$ & 170 & $0\,R\,183$ & $1\,R\,145$\\
74 & $0\,R\,47$ & $1\,R\,73$ & 171 & $0\,R\,158$ & $1\,R\,130$\\
75 & $0\,R\,88$ & $1\,R\,97$ & 172 & $0\,R\,139$ & $1\,R\,165$\\
76 & $0\,R\,47$ & $1\,R\,98$ & 173 & $0\,R\,158$ & $1\,R\,184$\\
77 & $1\,R\,77$ & $1\,L\,15$ & 174 & $0\,R\,185$ & $1\,R\,102$\\
78 & $0\,R\,90$ & $1\,R\,41$ & 175 & $0\,R\,31$ & $1\,R\,49$\\
79 & $0\,R\,99$ & $1\,R\,100$ & 176 & $0\,R\,186$ & $1\,R\,175$\\
80 & $0\,R\,41$ & $1\,R\,101$ & 177 & $0\,R\,134$ & $1\,R\,142$\\
81 & $0\,R\,102$ & $1\,R\,85$ & 178 & $0\,R\,32$ & $1\,R\,187$\\
82 & $0\,R\,103$ & $0\,L\,67$ & 179 & $0\,R\,144$ & $1\,R\,188$\\
83 & $0\,R\,104$ & $1\,R\,83$ & 180 & $0\,R\,189$ & $1\,R\,190$\\
84 & $0\,R\,105$ & $1\,R\,84$ & 181 & $0\,R\,40$ & $1\,R\,191$\\
85 & $0\,R\,106$ & $0\,L\,14$ & 182 & $0\,R\,168$ & $1\,R\,40$\\
86 & $0\,R\,62$ & $0\,R\,104$ & 183 & $0\,R\,192$ & $1\,R\,46$\\
87 & $0\,R\,107$ & $0\,L\,85$ & 184 & $0\,R\,153$ & $1\,R\,146$\\
88 & $0\,R\,30$ & $1\,R\,57$ & 185 & $0\,R\,165$ & $1\,R\,180$\\
89 & $0\,R\,108$ & $1\,R\,89$ & 186 & $0\,R\,119$ & $1\,R\,109$\\
90 & $0\,L\,60$ & $0\,R\,84$ & 187 & $0\,R\,193$ & $1\,R\,159$\\
91 & $0\,R\,90$ & $1\,L\,109$ & 188 & $0\,R\,169$ & $1\,R\,154$\\
92 & $0\,R\,97$ & $0\,L\,14$ & 189 & $0\,R\,158$ & $1\,R\,192$\\
93 & $0\,R\,64$ & $1\,R\,110$ & 190 & $0\,R\,179$ & $1\,R\,182$\\
94 & $0\,R\,109$ & $1\,R\,71$ & 191 & $0\,R\,182$ & $1\,R\,40$\\
95 & $0\,L\,60$ & $0\,R\,105$ & 192 & $0\,R\,149$ & $1\,R\,168$\\
96 & $0\,L\,59$ & $1\,R\,111$ & 193 & $0\,R\,72$ & $1\,R\,144$\\
\end{longtable}
\normalsize


\clearpage
\section{Complete dense PA194D transition table}
\label{app:table_d}
The second table uses the same entry notation and machine convention as
Appendix~\ref{app:table}. Its initial state is 0 and terminal label is 194.
All 388 rows agree with \file{logical_emitter/machine/pa194d.tm}.
\small
\begin{longtable}{@{}rll@{\hspace{1.3cm}}rll@{}}
\toprule
State & Read 0 & Read 1 & State & Read 0 & Read 1\\
\midrule\endfirsthead
\toprule
State & Read 0 & Read 1 & State & Read 0 & Read 1\\
\midrule\endhead
\midrule\multicolumn{6}{r}{Continued on the next page}\\\endfoot
\bottomrule\endlastfoot
0 & $1\,R\,1$ & $1\,R\,2$ & 97 & $0\,R\,36$ & $1\,R\,112$\\
1 & $1\,R\,3$ & $1\,R\,1$ & 98 & $0\,R\,111$ & $1\,R\,75$\\
2 & $0\,R\,0$ & $1\,R\,2$ & 99 & $0\,L\,63$ & $0\,R\,109$\\
3 & $0\,R\,4$ & $0\,R\,1$ & 100 & $0\,L\,62$ & $1\,R\,113$\\
4 & $1\,R\,5$ & $1\,L\,6$ & 101 & $0\,L\,90$ & $1\,R\,114$\\
5 & $1\,L\,7$ & $1\,L\,8$ & 102 & $0\,R\,115$ & $0\,L\,14$\\
6 & $0\,L\,7$ & $1\,L\,6$ & 103 & $0\,R\,69$ & $1\,R\,103$\\
7 & $0\,L\,9$ & $1\,L\,6$ & 104 & $0\,R\,116$ & $1\,R\,95$\\
8 & $0\,L\,10$ & $0\,L\,11$ & 105 & $0\,R\,116$ & $1\,R\,67$\\
9 & $1\,R\,12$ & $0\,L\,9$ & 106 & $0\,R\,117$ & $0\,L\,14$\\
10 & $0\,L\,9$ & $0\,L\,10$ & 107 & $0\,R\,118$ & $1\,L\,71$\\
11 & $1\,L\,8$ & $1\,L\,11$ & 108 & $0\,R\,119$ & $1\,R\,120$\\
12 & $0\,R\,12$ & $0\,R\,13$ & 109 & $0\,R\,14$ & $1\,R\,109$\\
13 & $0\,L\,14$ & $0\,R\,13$ & 110 & $0\,R\,66$ & $1\,R\,110$\\
14 & $1\,L\,15$ & $0\,L\,15$ & 111 & $0\,R\,121$ & $1\,L\,35$\\
15 & $0\,L\,15$ & $1\,L\,16$ & 112 & $0\,R\,122$ & $1\,R\,123$\\
16 & $0\,L\,17$ & $1\,L\,16$ & 113 & $0\,R\,88$ & $0\,L\,14$\\
17 & $0\,L\,18$ & $1\,L\,16$ & 114 & $0\,R\,90$ & $1\,L\,95$\\
18 & $0\,L\,19$ & $1\,L\,16$ & 115 & $0\,R\,124$ & $1\,R\,124$\\
19 & $0\,L\,20$ & $1\,L\,16$ & 116 & $0\,L\,63$ & $0\,R\,103$\\
20 & $0\,L\,21$ & $1\,L\,16$ & 117 & $0\,R\,68$ & $0\,R\,91$\\
21 & $0\,L\,22$ & $1\,L\,16$ & 118 & $0\,R\,58$ & $0\,R\,91$\\
22 & $0\,L\,23$ & $1\,L\,16$ & 119 & $0\,L\,86$ & $1\,R\,125$\\
23 & $0\,R\,23$ & $1\,R\,24$ & 120 & $0\,R\,126$ & $1\,R\,127$\\
24 & $0\,R\,25$ & $1\,R\,26$ & 121 & $0\,R\,58$ & $0\,R\,110$\\
25 & $0\,R\,27$ & $1\,R\,28$ & 122 & $0\,R\,99$ & $1\,R\,113$\\
26 & $0\,R\,29$ & $1\,R\,30$ & 123 & $0\,R\,79$ & $1\,R\,78$\\
27 & $0\,R\,31$ & $1\,R\,194$ & 124 & $0\,R\,88$ & $1\,R\,88$\\
28 & $0\,L\,32$ & $1\,R\,33$ & 125 & $0\,R\,128$ & $1\,R\,129$\\
29 & $0\,R\,34$ & $1\,R\,35$ & 126 & $0\,R\,32$ & $1\,R\,36$\\
30 & $0\,L\,36$ & $1\,L\,37$ & 127 & $0\,R\,31$ & $1\,R\,111$\\
31 & $0\,R\,38$ & $1\,R\,39$ & 128 & $0\,R\,130$ & $1\,R\,131$\\
32 & $0\,R\,40$ & $1\,R\,41$ & 129 & $0\,R\,132$ & $1\,R\,133$\\
33 & $0\,R\,42$ & $1\,R\,43$ & 130 & $0\,R\,134$ & $1\,R\,135$\\
34 & $0\,R\,44$ & $1\,R\,45$ & 131 & $0\,R\,136$ & $1\,R\,137$\\
35 & $0\,R\,46$ & $1\,R\,33$ & 132 & $0\,R\,88$ & $1\,R\,138$\\
36 & $0\,R\,47$ & $1\,R\,48$ & 133 & $0\,R\,139$ & $1\,R\,140$\\
37 & $0\,R\,49$ & $1\,R\,50$ & 134 & $0\,R\,97$ & $1\,R\,35$\\
38 & $0\,R\,51$ & $1\,R\,42$ & 135 & $0\,R\,141$ & $1\,R\,124$\\
39 & $0\,R\,52$ & $1\,R\,53$ & 136 & $0\,R\,142$ & $1\,R\,143$\\
40 & $0\,R\,54$ & $1\,R\,55$ & 137 & $0\,R\,144$ & $1\,R\,145$\\
41 & $0\,R\,56$ & $1\,R\,57$ & 138 & $0\,L\,137$ & $1\,R\,146$\\
42 & $0\,R\,58$ & $0\,R\,59$ & 139 & $0\,L\,147$ & $1\,R\,41$\\
43 & $0\,R\,60$ & $0\,L\,14$ & 140 & $0\,R\,148$ & $1\,R\,149$\\
44 & $0\,R\,61$ & $1\,R\,62$ & 141 & $0\,L\,148$ & $1\,R\,41$\\
45 & $0\,R\,63$ & $1\,R\,52$ & 142 & $0\,R\,28$ & $1\,R\,150$\\
46 & $0\,R\,64$ & $1\,R\,65$ & 143 & $0\,R\,151$ & $1\,R\,152$\\
47 & $0\,R\,58$ & $0\,R\,66$ & 144 & $0\,R\,153$ & $1\,R\,154$\\
48 & $0\,R\,67$ & $0\,L\,14$ & 145 & $0\,R\,155$ & $1\,R\,156$\\
49 & $0\,R\,68$ & $0\,R\,69$ & 146 & $0\,R\,135$ & $1\,R\,130$\\
50 & $0\,R\,70$ & $0\,L\,48$ & 147 & $0\,R\,157$ & $1\,R\,158$\\
51 & $0\,R\,64$ & $1\,R\,53$ & 148 & $0\,R\,159$ & $1\,R\,160$\\
52 & $0\,L\,71$ & $0\,L\,15$ & 149 & $0\,R\,161$ & $1\,R\,121$\\
53 & $0\,R\,72$ & $0\,L\,10$ & 150 & $0\,R\,162$ & $1\,R\,163$\\
54 & $0\,R\,73$ & $1\,R\,74$ & 151 & $0\,R\,155$ & $1\,R\,164$\\
55 & $0\,R\,75$ & $1\,R\,76$ & 152 & $0\,R\,138$ & $1\,R\,165$\\
56 & $0\,R\,77$ & $1\,R\,78$ & 153 & $0\,R\,134$ & $1\,R\,166$\\
57 & $0\,R\,79$ & $1\,R\,80$ & 154 & $0\,R\,167$ & $1\,R\,168$\\
58 & $0\,L\,81$ & $0\,R\,1$ & 155 & $0\,L\,169$ & $1\,R\,139$\\
59 & $0\,R\,82$ & $1\,R\,59$ & 156 & $0\,R\,170$ & $1\,R\,171$\\
60 & $0\,L\,73$ & $1\,L\,46$ & 157 & $0\,R\,41$ & $1\,R\,158$\\
61 & $0\,R\,83$ & $1\,R\,84$ & 158 & $0\,L\,172$ & $1\,L\,156$\\
62 & $0\,R\,54$ & $1\,R\,85$ & 159 & $0\,R\,76$ & $1\,R\,173$\\
63 & $0\,L\,86$ & $1\,R\,87$ & 160 & $0\,R\,158$ & $1\,R\,107$\\
64 & $0\,R\,58$ & $0\,R\,82$ & 161 & $0\,R\,174$ & $1\,R\,175$\\
65 & $0\,R\,88$ & $1\,R\,30$ & 162 & $0\,R\,134$ & $1\,R\,176$\\
66 & $0\,R\,89$ & $1\,R\,66$ & 163 & $0\,R\,76$ & $1\,L\,177$\\
67 & $0\,L\,90$ & $1\,L\,33$ & 164 & $0\,R\,172$ & $1\,R\,177$\\
68 & $0\,L\,81$ & $0\,R\,2$ & 165 & $0\,R\,166$ & $1\,R\,157$\\
69 & $0\,R\,91$ & $1\,R\,69$ & 166 & $0\,R\,168$ & $1\,R\,124$\\
70 & $0\,R\,72$ & $1\,R\,92$ & 167 & $0\,L\,129$ & $1\,R\,171$\\
71 & $0\,R\,46$ & $1\,R\,93$ & 168 & $0\,L\,151$ & $1\,L\,135$\\
72 & $0\,R\,42$ & $0\,L\,13$ & 169 & $0\,R\,178$ & $1\,R\,179$\\
73 & $0\,R\,94$ & $1\,R\,37$ & 170 & $0\,R\,180$ & $1\,R\,181$\\
74 & $0\,R\,30$ & $1\,R\,95$ & 171 & $0\,R\,182$ & $1\,R\,183$\\
75 & $0\,R\,42$ & $1\,R\,96$ & 172 & $0\,R\,139$ & $1\,R\,130$\\
76 & $0\,R\,97$ & $1\,R\,98$ & 173 & $0\,R\,184$ & $1\,R\,185$\\
77 & $0\,R\,99$ & $1\,L\,100$ & 174 & $0\,R\,31$ & $1\,R\,52$\\
78 & $0\,R\,30$ & $1\,R\,77$ & 175 & $0\,R\,186$ & $1\,R\,174$\\
79 & $0\,R\,93$ & $1\,R\,101$ & 176 & $0\,R\,141$ & $1\,R\,32$\\
80 & $0\,R\,30$ & $1\,R\,102$ & 177 & $0\,R\,179$ & $1\,R\,155$\\
81 & $1\,R\,81$ & $1\,L\,15$ & 178 & $0\,R\,182$ & $1\,R\,141$\\
82 & $0\,R\,103$ & $1\,R\,82$ & 179 & $0\,R\,181$ & $1\,R\,135$\\
83 & $0\,R\,94$ & $1\,R\,42$ & 180 & $0\,R\,154$ & $1\,R\,179$\\
84 & $0\,R\,104$ & $1\,R\,105$ & 181 & $0\,R\,183$ & $1\,R\,41$\\
85 & $0\,R\,42$ & $1\,R\,106$ & 182 & $0\,R\,97$ & $1\,R\,187$\\
86 & $0\,R\,107$ & $1\,R\,48$ & 183 & $0\,L\,130$ & $1\,R\,88$\\
87 & $0\,R\,108$ & $0\,L\,72$ & 184 & $0\,R\,163$ & $1\,R\,188$\\
88 & $0\,R\,68$ & $0\,R\,66$ & 185 & $0\,R\,189$ & $1\,R\,190$\\
89 & $0\,R\,109$ & $1\,R\,89$ & 186 & $0\,R\,92$ & $1\,R\,111$\\
90 & $0\,R\,49$ & $1\,R\,100$ & 187 & $0\,R\,107$ & $1\,R\,85$\\
91 & $0\,R\,110$ & $1\,R\,91$ & 188 & $0\,R\,147$ & $1\,R\,165$\\
92 & $0\,R\,30$ & $1\,R\,48$ & 189 & $0\,R\,191$ & $1\,R\,167$\\
93 & $0\,R\,47$ & $1\,R\,60$ & 190 & $0\,R\,169$ & $1\,R\,192$\\
94 & $0\,L\,63$ & $0\,R\,89$ & 191 & $0\,R\,182$ & $1\,R\,171$\\
95 & $0\,R\,94$ & $1\,L\,111$ & 192 & $0\,R\,193$ & $1\,R\,166$\\
96 & $0\,R\,101$ & $0\,L\,14$ & 193 & $0\,R\,191$ & $1\,R\,141$\\
\end{longtable}
\normalsize


\clearpage
\section{All 41 research packets}
\label{app:ledger}
The following ledger records the principal contribution of each packet.
Historical semantic claims are attributed to their supplied reports;
literal table identities and the present frontier comparison were freshly
checked. Full locators, archive identities and detailed evidence statuses
are in the CSV/JSON ledger.
\footnotesize
\begin{longtable}{@{}p{0.39\textwidth}p{0.57\textwidth}@{}}
\toprule
Packet & Principal contribution or outcome\\\midrule\endfirsthead
\toprule
Packet & Principal contribution or outcome\\\midrule\endhead
\midrule\multicolumn{2}{r}{Continued on the next page}\\\endfoot
\bottomrule\endlastfoot
\file{Below_U208_Research_2026-10-29.zip} & Full-frame windows, offset quotients and guarded pruning; source 3,909$\to$3,904; no accepted below-208 table.\\[2pt]
\file{ConPA_below208_research_2026-10-29.zip} & Finite-model role-collapse objection, 279,936 fixed-association parameter cases, fold synthesis; no sub-208.\\[2pt]
\file{PA195_ConPA_research_2026-11-14.zip} & Positive 195S and 195C: unchanged separate 190-token AX word, 67$-$6 readers; full source/runtime/startup certificates plus independent literal-tape audits.\\[2pt]
\file{PA_consistency_TM_research_2026-11-13.zip} & 159-node five-zero source and 157-node combined representation; exact source CFG checks, no smaller physical allocation.\\[2pt]
\file{U208_ConPA_2026-10-29.zip} & Particular native controls 172/211 history fold reduces a U209 parent to U208.\\[2pt]
\file{U208_ConPA_Conditional_Context_2026-10-29.zip} & Conditional-entry proof; same two original U208 graph classes.\\[2pt]
\file{U208_ConPA_Entry_Context_2026-10-29.zip} & Full entry-context proof; same two original classes.\\[2pt]
\file{U208_ConPA_History_2026-10-29.zip} & Finite-history proof with four exceptional blocks; same original saving.\\[2pt]
\file{conpa_tm_comparison_and_reduction_search.zip} & Experimental MP-result-in-H and MP-without-zero-push 196 variants; rejected Shift layouts halt at 43,563.\\[2pt]
\file{conpa_tm_comparison_and_search.zip} & A/D register-swap 196; rejected two-rotation 195 fails 93 contracts.\\[2pt]
\file{conpa_tm_study_2026-11-13.zip} & Combined196 = H-only Shift + MP-result-in-H: 3,788 sites, 4,928 contracts, 57,922 startup; direct-negation 197 too.\\[2pt]
\file{pa194C_conpa_2026-11-13 (2).zip} & PA194C; same literal cached-head machine as other C packets. Broad comparison and below194 searches.\\[2pt]
\file{pa194_conpa_2026-11-13 (2).zip} & Different separate 190-token association: 68 readers and 8 simultaneous folds yield a third positive194 graph. Also contains a distinct intermediate195 separate 190-token layout with 67 readers and6 folds. The actual intermediate source is3673 sites/4780 contracts despite a clerical report-table diagonal/3700 label.\\[2pt]
\file{pa194_conpa_2026-11-13.zip} & Three-way saving composition: separate 190-token emitter, 67 raw AX readers and 7 simultaneous folds yield primary PA194; the A/D register permutation gives a second nonisomorphic194. Both include complete 4780 symbolic contracts, exact startup, and a written PA calibration.\\[2pt]
\file{pa194_new_designs_2026-11-13.zip} & PA194C plus PA194CS. CS uses equality symmetry to remove three source sites, adds a reader and folds41 into330, remaining194.\\[2pt]
\file{pa194c_conpa.zip} & PA194C; also eight-fold cached variant pa194altC as checked search input, not principal theorem.\\[2pt]
\file{pa194c_conpa_2026-11-13 (1).zip} & PA194C plus PA194CZ zero-register-jump variant at source 26, justified by263244 positivity-mask pairs.\\[2pt]
\file{pa194c_conpa_2026-11-13.zip} & PA194C; same dense/native literal machine as four other C packets; alternate independent audit and structural search records.\\[2pt]
\file{pa195_conpa_2026-11-13.zip} & Positive original PA195: unchanged diagonal191 program, 68$-$7 reader allocation, detailed ordinary PA calibration.\\[2pt]
\file{pa195_diagonal_combined_2026-11-13.zip} & PA195DC combines bundled diagonal191 and7 folds with cached-head Shift/MP; 3676 sites,4784 contracts,57922 startup, six changed native rows. Optional PA195R and PA195N implement parameter-permuted and negative-vacuity emitters at195, with explicitly limited endpoint claims.\\[2pt]
\file{pa195_followup_research.zip} & M195S8 keeps diagonal191 but pays one extra scanner control and obtains8 folds:135+68-8=195. It rejects the associated194 quotient at source ordinal1003 because seven address zeros trigger the old scanner. Includes additional equivalent source-calculus families and bounded searches.\\[2pt]
\file{pa196_research_2026-11-13.zip} & Exact 157-node representation of 3,812-site original program; loses contextual sharing, no physical reduction.\\[2pt]
\file{pa208_research_2026-10-29.zip} & Eighteen further helper layouts; 27 designated U208 presentations give 20 graphs, not two.\\[2pt]
\file{pa_below196_research.zip} & Eight equivalent gates, membership/equality orientations, broad grammar searches; original201-token best 196.\\[2pt]
\file{pa_consistency_comparison_2026_11_13.zip} & Modal-head/explosion split, twelve semantic compositions, joint gate; physical results 196–200.\\[2pt]
\file{pa_consistency_comparison_and_new196_2026-11-13.zip} & Diagonal196 (3,700 sites), Separate196 (3,673), Scanner-best196; principal logical/source parents of positive195.\\[2pt]
\file{pa_consistency_cross_archive.zip} & Four new196 combinations: diagonal/separate AX + Combined controller, each optionally zero-jump; sites 3,676/3,649 and startups57,922/58,168.\\[2pt]
\file{pa_consistency_reduction.zip} & Rebracketed 196, H-only 196, integrated gate 203/206; two failed 195 diagnostics.\\[2pt]
\file{pa_consistency_reduction_audit_2026-11-14.zip} & Four further experimental MP-H compositions, two 196-state and two 197-state; parameter/order/association/helper/scanner negative searches.\\[2pt]
\file{pa_consistency_research.zip} & Two exact 158-node sources and repaired 196; nominal 195 needs added scanner state.\\[2pt]
\file{pa_consistency_research_2026-10-29.zip} & Inductive 3,909$\to$3,904 specialization; no physical sub-208.\\[2pt]
\file{pa_consistency_research_2026-11-13 (2).zip} & Twelve-rotation 196, diagonal 198, successful-MP halting-source experiment; no sub-196.\\[2pt]
\file{pa_consistency_research_2026-11-13.zip} & Earlier diagonal-vacuity 199 and short-negation 197.\\[2pt]
\file{pa_consistency_research_2026-11-14.zip} & NEG18: `INC7; INC7; node64`, 18-site encoded negation; complete source 3,441 and 197 states (134+68$-$5); independent macro and scratch-zero proofs.\\[2pt]
\file{pa_tm_comparison_and_reductions.zip} & Two rotated 196-state graphs; distributed-gate 198 with 3,773 source sites.\\[2pt]
\file{pa_tm_extended_research_2026-11-13.zip} & Four closing-decrement variants and stronger zero invariants; all change the same row, hence not stackable.\\[2pt]
\file{pa_tm_research_2026-11-13 (2).zip} & Zerojump 196 (4,959 contracts, 58,402 startup); parity/disequality experiments; rejected 195 fails source 26.\\[2pt]
\file{pa_tm_research_2026-11-13.zip} & Shorter-path 196 and independent physical checking; detailed direct PA reverse-translation strategy.\\[2pt]
\file{pa_unary_power_cover.zip} & Unary weak-covering theory and PA proof translations; no new compiled small table.\\[2pt]
\file{u208_comparison_research_2026-10-29.zip} & Original two-class audit; weak/exact schemas; rejected 207.\\[2pt]
\file{u208_research_2026-10-29.zip} & Primary/alternative U208 and symbolic contracts; same initial history fold.\\[2pt]
\end{longtable}
\normalsize

\end{document}
